\lstdefinelanguage{B}{
  morekeywords={main, auto, while, if, else, return},
  sensitive=true,
  morecomment=[l]{//},
  morecomment=[s]{/*}{*/},
  morestring=[b]",
}

\chapter{Технологический раздел}

\section{Сгенерированные классы анализаторов}

В результате работы ANTLR будут сгенерированы следующие классы:
\begin{itemize}
	\item bLexer — лексический анализатор.
	\item bParser — синтаксический анализатор.
	\item bVisitor — абстрактный класс посетителя для обхода дерева.
	\item bListener — абстрактный класс слушателя, с пустыми методами.
    \item bBaseVisitor - реализация интерфейса bVisitor c пустыми переопределенными методами.
    \item bBaseListener - реализация интерфейса bListener с пустыми переопределенными методами.
\end{itemize}

Для обхода дерева используется паттерн Visitor: создаётся класс-наследник от bVisitor, в котором переопределяются методы visit<RuleName>. 
Каждый из этих методов вызывается вручную при обходе дерева, а возврат значений и переход к дочерним узлам осуществляется через вызов visit(ctx->child) или visitChildren(ctx). 
Это позволяет точно управлять порядком обхода и обработкой содержимого узлов дерева, соответствующих правилам исходной грамматики.

\section{Обнаружение ошибок}

Все ошибки, обнаруженные на этапах лексического и синтаксического анализов, обрабатываются средствами ANTLR.
Возникающие при этом исключения обрабатываются и выводится текст соответствующей ошибки.
Также при построении графа вызовов обрабатываются ошибки вызова функции до ее объявления.

\section{Пример компиляции программы}

Пример программы, написанный на рассматриваемом языке B представлен на листинге \ref{lst:while-b}.
Данная программа реализует сортировку пузырьком.

\begin{lstlisting}[
	caption={Программа на рассматриваемом языкe B},
	label={lst:while-b},
	language={B},
	gobble=8
]
        main() {
            auto i;
            i = 0;
            while (i <= 10) {
                print(i);
                i++;
                i++;
            }

            auto j;
            j = 0;
            while (j <= 10) {
                print(j);
                j++;
            }

            auto k;
            k = 0;
        }
\end{lstlisting}

В листинге в приложении В приведен сгенерированный LLVM IR код для представленной на листинге программы.